\section{定積分の部分積分}\label{節:定積分の部分積分}

定積分を求めるとき,不定積分を求めてから上下端を代入してもよい.しかし最初から定積分の公式
\[
 \int_a^b u(x)v'(x)\dd x=[u(x)v(x)]_a^b-\int_a^b u'(x)v(x)\dd x
\]
を使えば,端点で消える項がすぐに分かる.本節では,定積分を端点の項 $[uv]_a^b$ と,残りの定積分 $\int_a^b u'v\dd x$ の2つに分けて計算する.

端点の項は
\[
 [u(x)v(x)]_a^b=u(b)v(b)-u(a)v(a)
\]
であり,数を代入するだけで値が決まる.特に,端点で $u$ か $v$ のどちらかが $0$ になれば,その端点の寄与はなくなる.したがって,残る積分が簡単になるかに加えて,\,$u(a),u(b),v(a),v(b)$ の値も計算の見通しに関わる.端点の項を先に計算しておけば,あとは残りの定積分に集中できる.

\begin{bookpoint}
\textbf{POINT}\par
端点の項 $[uv]_a^b$ を先に計算し,残る $\int_a^b u'v\dd x$ を求めて引く.
\end{bookpoint}

\begin{example}{端点の値}
$\int_0^{\frac\pi2}x\sin x\dd x$ を求めよ.
\end{example}

\paragraphbf{方針}
$x$ を微分し,\,$\sin x$ を積分する.端点の項は上端でも下端でも $0$ になる.

\paragraphbf{解答}
$1$

\paragraphbf{解法}
$u=x$,\,$v'=\sin x$ と選ぶと,\,$u'=1$,\,$v=-\cos x$ なので,
\[
\begin{aligned}
 \int_0^{\frac\pi2}x\sin x\dd x
 &=[-x\cos x]_0^{\frac\pi2}
   +\int_0^{\frac\pi2}\cos x\dd x\\
 &=0+[\sin x]_0^{\frac\pi2}\\
 &=1
\end{aligned}
\]
となる.上端では $\cos\frac\pi2=0$,下端では $x=0$ のため,端点の項はともに $0$ になった.2つの端点で $0$ になる理由は異なり,一方だけを見て項全体を消しているわけではない.

\paragraphbf{値の符号による確認}
積分される関数 $x\sin x$ はこの区間で非負だから,求めた値が正であることは自然である.仮に $-1$ になったら,\,$\sin x$ の原始関数の負号か,部分積分の前の負号を確かめるべきだと分かる.

\begin{example}{部分積分と割り算}
$\int_0^1x\log(1+x)\dd x$ を求めよ.
\end{example}

\paragraphbf{方針}
$\log(1+x)$ を微分し,\,$x$ を積分する.残る有理関数の積分は,分子を分母で割ってから計算する.

\paragraphbf{解答}
$\frac14$

\paragraphbf{解法}
$\log(1+x)$ を微分すると $\frac1{1+x}$ になる.\,$x$ の原始関数は $\frac{x^2}{2}$ だから,残る積分は有理関数になる.そこで
\[
 u=\log(1+x),\qquad v'=x
\]
と選ぶと,
\[
 \int_0^1x\log(1+x)\dd x
 =\left[\frac{x^2}{2}\log(1+x)\right]_0^1
  -\frac12\int_0^1\frac{x^2}{1+x}\dd x
\]
を得る.第1項は $\frac12\log2$ である.第2項では,分子の次数が分母より高いので,\chapterref{章:基本積分}で学んだ割り算を使う.
\[
 x^2=(1+x)(x-1)+1
\]
より,
\[
 \frac{x^2}{1+x}=x-1+\frac1{1+x}
\]
である.したがって
\[
\begin{aligned}
 \int_0^1\frac{x^2}{1+x}\dd x
 &=\left[\frac{x^2}{2}-x+\log(1+x)\right]_0^1\\
 &=-\frac12+\log2
\end{aligned}
\]
となる.これを戻すと
\[
\begin{aligned}
 \int_0^1x\log(1+x)\dd x
 &=\frac12\log2-\frac12\left(-\frac12+\log2\right)\\
 &=\frac14
\end{aligned}
\]
を得る.

\paragraphbf{消えた対数}
部分積分の直後には端点の項に対数が残り,残りの積分にも対数が現れたが,最後に差を取ると打ち消し合った.途中の式の見た目だけで,答えにも対数が必ず残ると決めることはできないのである.
